\chapter{Hypoglycemia Symptoms}

\section*{Definition}
Hypoglycemia is a clinical syndrome of low plasma glucose accompanied by characteristic autonomic and neuroglycopenic symptoms, defined biochemically as serum glucose \ensuremath{\leq}70 mg/dL (3.9 mmol/L) in the general population. Whipple's triad confirms the diagnosis: symptoms of hypoglycemia, low plasma glucose, and symptom resolution after glucose administration.

\section*{What Happens in Your Body}
As glucose declines below 70 mg/dL, the sympathetic nervous system and adrenal medulla release epinephrine, norepinephrine, and acetylcholine, producing autonomic warning symptoms (tremor, palpitations, sweating). At levels <54 mg/dL (3.0 mmol/L), neuroglycopenic symptoms emerge from cerebral glucose deprivation: confusion, seizure, and coma. CNS glucose supply is critically dependent on plasma glucose because neurons cannot synthesize or store significant glucose, and alternative fuels (ketones) are insufficient under most conditions.

\section*{Causes}
\begin{itemize}
  \item \textbf{Diabetes-related:} Excess insulin (exogenous, sulfonylureas, meglitinides), missed meals, exercise without carbohydrate adjustment, alcohol consumption, renal impairment delaying insulin clearance
  \item \textbf{Non-diabetic fasting:} Insulinoma (endogenous hyperinsulinism), insulin autoimmune syndrome, non-islet cell tumor hypoglycemia (IGF-2), adrenal insufficiency, liver failure, renal failure, sepsis
  \item \textbf{Non-diabetic reactive:} Post-gastric bypass dumping syndrome, idiopathic postprandial syndrome, early type 2 diabetes (inappropriate insulin secretion)
  \item \textbf{Factitious:} Surreptitious insulin or sulfonylurea use (healthcare workers, caregivers)
  \item \textbf{Drug-induced:} Alcohol (impaired gluconeogenesis), salicylates, pentamidine, quinolones, beta-blockers (mask symptoms)
\end{itemize}

\section*{When to See a Doctor}
See your doctor if:
\begin{itemize}
  \item Severe hypoglycemia requiring assistance from another person (glucagon, EMS)
  \item Recurrent unexplained episodes despite adjustments in diabetes management
  \item Loss of hypoglycemia awareness (hypoglycemia unawareness, common in longstanding type 1 diabetes)
  \item Fasting or post-prandial hypoglycemia in non-diabetic people
  \item Altered mental status, seizure, or loss of consciousness
\end{itemize}

\section*{Related Symptoms}
\begin{itemize}
  \item Diaphoresis, tremor, palpitations, hunger, anxiety, confusion, weakness, seizure, coma
\end{itemize}
\textbf{Important:} Recurrent hypoglycemia can lead to hypoglycemia-associated autonomic failure (HAAF), where the glucose threshold for autonomic activation shifts to lower levels, diminishing or eliminating warning symptoms.

\section*{Self-Care / Home Management}
\begin{itemize}
  \item Follow the Rule of 15: If glucose <70 mg/dL, consume 15 g rapid-acting carbohydrate, recheck in 15 minutes, repeat if still <70
  \item Rapid-acting sources: 4 oz (120 mL) juice or regular soda, 1 tablespoon sugar or honey, 3--4 glucose tablets
  \item Follow with a snack containing protein and complex carbohydrate after stabilization
\end{itemize}

\section*{How Your Doctor Will Evaluate You}
\begin{itemize}
  \item Document Whipple's triad with home glucose monitoring or continuous glucose monitoring (CGM) data
  \item Medication review for diabetics: timing, dosing, and concordance
  \item Fasting (72-hour) supervised test with plasma glucose, insulin, C-peptide, proinsulin, beta-hydroxybutyrate, and sulfonylurea screen
  \item Imaging for insulinoma: CT, MRI, endoscopic ultrasound, or octreotide scan
\end{itemize}

\section*{Treatment Options}
\begin{itemize}
  \item Acute: Oral glucose if conscious; IM glucagon (1 mg adult/0.5 mg child) or IV dextrose 50\% (25 g) for severe hypoglycemia with altered mental status
  \item Diabetes management adjustment: Reduce or re-time insulin/oral secretagogues, increase carbohydrate intake before exercise, utilize CGM with low-glucose alerts
  \item Insulinoma: Surgical enucleation or resection; medical management (diazoxide, octreotide, everolimus) for inoperable cases
  \item Post-bariatric: Dietary modification (low-glycemic, small frequent meals), acarbose, GLP-1 receptor agonist adjustments
  \item Hormonal replacement for adrenal insufficiency or growth hormone deficiency
\end{itemize}

\section*{References}

\begin{thebibliography}{99}
\bibitem{harrison-hypogly} Longo DL, et al. \emph{Harrison's Principles of Internal Medicine}. 21st ed. McGraw-Hill; 2022. Chapter 411: Hypoglycemia.
\bibitem{uptodate-hypogly} Cryer PE, et al. Hypoglycemia in adults with diabetes mellitus. In: UpToDate, Post TW (Ed), UpToDate, Waltham, MA; 2025.
\bibitem{cryer} Cryer PE. Hypoglycemia: pathophysiology and management. \emph{N Engl J Med}. 2022;386(19):1814-1824.
\bibitem{ada-hypogly} American Diabetes Association. Standards of Care in Diabetes: Hypoglycemia. \emph{Diabetes Care}. 2024;47(Suppl 1):S134-S144.
\bibitem{williams-hypogly} Melmed S, et al. \emph{Williams Textbook of Endocrinology}. 15th ed. Elsevier; 2023. Chapter 35: Hypoglycemia.
\bibitem{seaquist} Seaquist ER, et al. Hypoglycemia and diabetes: a consensus report. \emph{Diabetes Care}. 2023;46(7):1389-1401.
\end{thebibliography}

\clearpage
